% TOP_PROBLEM: {"id":30007582,"problem_number":"LOCAL-30007582","title":"Risk preferences versus complexity errors","category_id":21,"category":{"id":21,"name":"economics","display_name":"Economics","description":"Open economic theory statements, published research agendas, and proposed scoped specifications; question provenance and historical priority are qualified per record.","slug":"economics","order_index":21,"created_at":"2026-10-08T00:00:00Z"},"status":"provenance_pending","external_url":null,"legacy_ids":[],"published":false,"statement":"Separate stable preferences over risk from complexity-induced decision errors using the specified variation in lottery representation and incentives.","economics":{"original_id":71,"field":"Behavioral and experimental economics","kind":"identification","formalization_basis":"proposed_specification","source_status":"not_verified","priority_status":"not_established","priority_note":"The cited experiment motivates this decomposition but no primary locator explicitly posing the consequential-field version was verified. No first-statement attribution is made.","earliest_verified_reference":null,"current_reference":{"authors":["Ryan Oprea"],"title":"Decisions under Risk Are Decisions under Complexity","year":2024,"url":"https://www.aeaweb.org/articles?id=10.1257/aer.20221227","doi":"10.1257/aer.20221227","locator":"Abstract; publisher metadata, pp. 3789–3811. Full paper was not accessible in this audit.","evidence_summary":"The abstract reports lottery-like anomalies in tasks without risk. It does not establish the broad consequential-field decomposition in this entry as an explicitly open problem."},"formalization_note":"This is an explicitly proposed measurement design. The additive error equation and exclusion restrictions are our specification, not a theorem or equation attributed to Oprea."},"statusQualification":"Provenance pending: an explicit published statement of this exact open question was not verified. This is a catalogue-authored candidate specification, not a certified open conjecture. This is an explicitly proposed measurement design. The additive error equation and exclusion restrictions are our specification, not a theorem or equation attributed to Oprea.","statusReviewedAt":"2026-10-08"}
% FIRST_POSTED: 2026-10-08
% ENTRY_KIND: statement_only
% !TeX program = lualatex
\documentclass[11pt]{article}
\usepackage{fontspec}
\usepackage[a4paper,margin=25mm]{geometry}
\usepackage{amsmath,amssymb,mathtools}
\usepackage{xurl}
\usepackage[hidelinks]{hyperref}
\newcommand{\sref}[2]{\hyperlink{source:#1:#2}{[S#2]}}
\setlength{\emergencystretch}{3em}
\begin{document}
% problemId:30007582
\section*{Risk preferences versus complexity errors}\label{30007582}
\subsection{Definitions and mathematical statement}
Consider individuals \(i\) choosing between a sure payment and a binary lottery \(L=(p,x,y)\), where \(p\in[0,1]\) is the probability of receiving \(x\), and \(y\) is received otherwise. Let \(C\in\{0,1\}\) randomly assign a simple or complex representation while preserving \(L\), payment rules, and information. Let \(a_i(L,C)\) be the reported certainty equivalent. A preference model specifies utility \(u_i\), probability weighting \(w_i\), and a reference point \(r_i\); its latent certainty equivalent \(v_i(L)\) satisfies \(u_i(v_i)=w_i(p)u_i(x)+(1-w_i(p))u_i(y)\). Observations obey \(a_i(L,C)=v_i(L)+e_i(L,C)\), where \(e_i\) is a decision or reporting error whose conditional mean need not vanish. The observable causal complexity effect is
\[\tau(L)=\mathbb E[a_i(L,1)-a_i(L,0)].\]
Identify which restrictions, auxiliary tasks, or repeated measurements permit separation of \(v_i\) from systematic \(e_i\), rather than attributing every pattern to preferences. Random assignment identifies \(\tau\) but does not identify the error-free preference distribution without further assumptions. A field extension should specify a target population \(P\), comparable lotteries, and consequential implementation; estimate the corresponding \(\tau_P(L)\) and the predictive performance of recovered preferences for held-out decisions. The unresolved empirical question concerns identification and transportability of this decomposition, including whether representations alter reference points or attention as well as computational burden.

\par\textbf{Formulation and scope.} This is an explicitly proposed measurement design. The additive error equation and exclusion restrictions are our specification, not a theorem or equation attributed to Oprea.

\par\textbf{Open-question provenance.} An explicit published open-question statement was not verified. Sources below are supporting literature and must not be treated as a first listing.

\par\textbf{Historical priority.} The cited experiment motivates this decomposition but no primary locator explicitly posing the consequential-field version was verified. No first-statement attribution is made.

\subsection{Short English statement}
Separate stable preferences over risk from complexity-induced decision errors using the specified variation in lottery representation and incentives.

\subsection{Sources}
\begin{itemize}\raggedright
\item[\textnormal{[S1]}]\hypertarget{source:30007582:1}{} Ryan Oprea. 2024. Decisions under Risk Are Decisions under Complexity. Abstract; publisher metadata, pp. 3789–3811. Full paper was not accessible in this audit..
\url{https://www.aeaweb.org/articles?id=10.1257/aer.20221227}.
DOI: 10.1257/aer.20221227.
{\small\textit{Supporting or current literature; see evidence qualification. The abstract reports lottery-like anomalies in tasks without risk. It does not establish the broad consequential-field decomposition in this entry as an explicitly open problem.}}
\end{itemize}
\end{document}
